\begin{resumo}[Abstract]

This work evaluates the probabilistic calibration of a Temporal Fusion Transformer applied to multi-horizon asset return forecasting and characterizes the relative contribution of indicator families under a reproducible and auditable experimental protocol. The research problem stems from the difficulty of fairly comparing model configurations in out-of-sample settings when evaluation is restricted to point metrics and when predictive uncertainty, temporal alignment, and execution traceability are not treated as explicit methodological requirements. The scope is single-asset (one model per asset), with AAPL as the initial pilot and main horizons $h+1$ and $h+7$. The methodology integrates market data, technical indicators, aggregated sentiment, and fundamentals, with versioned feature engineering, leakage control, and quantile-mode training of the Temporal Fusion Transformer ($p10$, $p50$, $p90$) under a post-guardrail variant policy as the primary probabilistic evaluation object. The experimental protocol uses explicit temporal partitioning, multiple seeds, cohort governance by \textit{parent\_sweep\_id}, strict alignment by target timestamp, and paired comparison against pre-declared baselines (random walk, AR(1), historical mean, EWMA-vol, historical empirical quantiles) via Diebold-Mariano with HAC correction and Holm-Bonferroni adjustment. Feature-family contribution analysis combines the Variable Selection Network, permutation importance, and explanatory ablation. Experimental artifacts are persisted in immutable persistence layers and in a derived, reconstructible analytical layer, enabling paired analysis without retraining. Confirmatory results and the contribution analysis will be consolidated under a versioned pre-registration.

\textbf{[Version note — to be removed in the final defense version]} This abstract is a living version, updated at each advising round as new evidence, methodological refinements, and consolidated confirmatory results from Phase B become available. The final version submitted to the examining committee will not include this note.

\textbf{Keywords}: Temporal Fusion Transformer. Quantile forecasting. Probabilistic calibration. Experimental reproducibility. Out-of-sample evaluation.

\end{resumo}
